\section{Introduction}

LLM-based systems can generate many plausible GPU kernels for one tensor
program. The system then selects a candidate using measured latency. This makes
evaluation part of the search algorithm: as the candidate budget grows, so does
the chance of selecting a favorable timing fluctuation. Reusing that measurement
as final evidence produces a systems analogue of test-set overfitting. A second
risk comes from the evaluator itself. Incorrect state initialization, model
construction inside timed regions, or permissive fallback paths can create a
large apparent speedup without a better kernel.

Existing kernel benchmarks emphasize correctness and performance over practical
baselines~\citep{kernelbench2025,zhang2026kernelbenchverified,lange2025robustkbench},
while recent harness work stresses constrained execution and artifact
archival~\citep{shui2026harness}. Systems measurement research has separately
shown that small experimental choices can reverse performance
conclusions~\citep{mytkowicz2009wrongdata,curtsinger2013stabilizer}. We connect
these concerns through a promotion protocol for generated kernels. The protocol
uses one dataset to screen candidates and fresh process-level data to confirm a
frozen winner.

We ask three questions with separate evidence paths. \textbf{RQ1:
repeatability.} Across a corrected KernelBench subset, how many valid tasks,
screening wins, and independently confirmed wins remain under a fixed
three-candidate budget? \textbf{RQ2: selection bias.} Across an easy
deterministic grid and a calibrated near-threshold fused8 stress test, how do
apparent and independently confirmed wins change with candidate budget?
\textbf{RQ3: evaluator validity.} In a
separate disposable-process ablation, how much apparent speedup is created by
reconstructing and transferring a reference model inside the measured call?

The contribution is an evaluation design rather than a generation model: a
contract-valid persistent lifecycle with frozen provenance; screen-then-confirm
timing across randomized paired blocks and fresh processes; fail-closed null,
known-slowdown, and lifecycle controls; and an artifact ledger that makes both
the candidate and evaluator auditable.

The claim ladder is deliberately strict:
\[
\text{contract-valid and correct}
\;\rightarrow\; \text{beats compiler}
\;\rightarrow\; \text{beats eager}
\;\rightarrow\; \text{independently confirmed}.
\]
Each arrow is a separate empirical claim.

Recent verifier studies show that kernel correctness requires more than a
single numerical comparison~\citep{shah2026contractgrade,li2026correctslow}.
Accordingly, our protocol treats correctness, output contracts, lifecycle
validity, and performance as distinct gates. The historical fused8 example is
motivating evidence only. The corrected campaign reported here passed its
prespecified control gates and is derived from checksummed raw artifacts.
